\documentclass[11pt,a4paper]{article}

% ---------- Packages ----------
\usepackage[margin=2.5cm]{geometry}
\usepackage{amsmath,amssymb,amsfonts}
\usepackage{graphicx}
\usepackage{booktabs}
\usepackage{siunitx}
\usepackage{hyperref}
\usepackage{url}
\usepackage{enumitem}
\usepackage{xcolor}
\usepackage{caption}
\usepackage{subcaption}
\usepackage{graphicx}
\usepackage{float}
\usepackage{caption}


\title{Deep Learning Project - Text Classification}
\author{Paul Brunner, Francesco Scivoletto, Moritz Zideck}
\date{\today}

\hypersetup{
    colorlinks=true,
    linkcolor=blue,
    citecolor=blue,
    urlcolor=blue,
    pdfauthor={Your Name},
    pdftitle={TITLE OF THE DEEP LEARNING PROJECT}
}

\begin{document}

\maketitle

\tableofcontents
\newpage

\section{Introduction}
    Modern GPT models have revolutionized natural language processing by enabling machines to understand and generate human-like text.
    A natural question that arises in this context is whether we can reliably determine the origin of a given text sample, i.e., whether it was written by a human or generated by a particular language model.
    This task is framed as a text classification problem.

    The goal of this project is to explore different approaches for classifying text samples according to their source, including various GPT-based models (e.g., Gemma, ChatGPT, LLaMA) and human-authored text, using a range of deep learning techniques and architectures.
     Since a labeled training dataset is not readily available, the first step is to construct such a dataset by scraping text samples from the web, which are then used to train a custom multi-head self-attention transformer model.

    Model performance is evaluated on a held-out dataset in order to assess its generalization capabilities and to identify potential areas for improvement.
    The evaluation samples primarily cover topics in science, technology, and medicine, which in turn requires the training set to be filtered and processed so that it contains sufficiently similar and relevant examples for a meaningful evaluation.

\section{Dataset}

Since no single dataset exists that covers all required model classes, a custom dataset was constructed by combining multiple publicly available sources. The goal was to obtain a balanced and diverse collection of human-written and machine-generated texts across several domains.

\subsection{Data Sources}

The dataset is composed of samples from the following sources:

\begin{itemize}
    \item \textbf{HC3 (Human ChatGPT Comparison Corpus)} \\
    Contains paired human-written and ChatGPT-generated answers.
    
    \item \textbf{OpenTuringBench (OTB)} \\
    Provides large-scale outputs from multiple LLMs such as Gemma and LLaMA.
    
    \item \textbf{GSINGH1 Dataset} \\
    Contains additional samples from multiple models and human-written texts.
\end{itemize}

Additionally, smaller datasets such as \texttt{claude\_multiround\_chat\_30k} were incorporated to increase diversity, especially for Anthropic-style outputs.

\subsection{Dataset Construction}

The dataset is generated using a unified pipeline that:

\begin{enumerate}
    \item Loads raw data from each source
    \item Normalizes and cleans text samples
    \item Maps model names to a common label space
    \item Filters samples by length
    \item Combines all sources into a single JSON dataset
\end{enumerate}

Only texts within a specific word range are kept to ensure consistency in style and difficulty of classification. In our experiments, we use:

\begin{itemize}
    \item Minimum length: 100 words
    \item Maximum length: 150 words
\end{itemize}

This filtering removes extremely short or excessively long samples that could bias the model.

\subsection{Final Dataset Statistics}

\begin{table}[h]
\centering
\begin{tabular}{lrrrrrr}
\hline
\textbf{Source} & \textbf{Total} & \textbf{Human} & \textbf{ChatGPT} & \textbf{Mistral} & \textbf{Gemma} & \textbf{LLaMA} \\
\hline
HC3      & 62,149  & 36,366 & 25,783 & 0      & 0      & 0      \\
GSINGH1  & 85,508  & 18,910 & 24,198 & 14,835 & 13,154 & 14,411 \\
OTB      & 187,760 & 0      & 0      & 64,596 & 54,710 & 68,454 \\
\hline
\textbf{Total} & \textbf{335,417} & \textbf{55,276} & \textbf{49,981} & \textbf{79,431} & \textbf{67,864} & \textbf{82,865} \\
\hline
\end{tabular}
\caption{Dataset distribution across sources and model classes}
\end{table}

\subsection{Dataset Format}

Each entry in the dataset follows a unified JSON structure:

\begin{verbatim}
{
  "model": "anthropic",
  "text": "...",
  "topic": "biology",
  "origin": "claude_multiround_chat_30k",
  "length": 135
}
\end{verbatim}

The fields are defined as follows:

\begin{itemize}
    \item \textbf{model}: source label (e.g., human, chatgpt, gemma, llama)
    \item \textbf{text}: the text sample
    \item \textbf{topic}: domain category (e.g., physics, biology, history)
    \item \textbf{origin}: dataset source
    \item \textbf{length}: number of words in the sample
\end{itemize}

\subsection{Dataset Characteristics}

The dataset is designed to capture key differences between human-written and AI-generated text, focusing on stylistic patterns and linguistic characteristics. It includes samples from multiple model families such as OpenAI, Meta, Google, and Anthropic, allowing the model to learn distinctions between different generation approaches. In addition, the dataset covers a range of domains, including science, technology, and medicine, ensuring that the classification task is not limited to a single topic area and generalizes across different types of content.

\section{Text preprocessing and data loading}
    Since the text samples are straight for each dataset and do not follow a common normalization pattern, they first have to be normalized.
 Normalization means that unwanted noise is removed while important stylistic information is preserved.
    For this purpose, all text samples undergo the following transformations:
    \begin{itemize}
    \label{normalisation}
    \item Unicode normalization: \verb|text = unicodedata.normalize("NFKC", text)|
    \item Normalize line endings: \verb|text = text.replace("\r\n", "\n").replace("\r", "\n")|
    \item Collapse whitespace: \verb|text = re.sub(r"[ \t]+", " ", text)|
    \item Collapse multiple newlines: \verb|text = re.sub(r"\n+", "\n", text)|
    \end{itemize}
    These help to standardize the text samples and reduce variability that could hinder model training.
    Since multible models use a vocabulary of words or subwords, the normalisation helps to ensure that the same word is represented consistently across the dataset, improving the model's ability to learn meaningful patterns.
    Next to normalisation the dataset is split into training, validation and test set.
    Furthermore the dataloader also alows to select the size of the dataset and to shuffle the data for training based on a deterministic seed.
    \section{Deep Neural Network Architecture}
    One of the model used in this project is a custom numpy implementation of a multi-head self attention transformer.
    It is based on the multi layer baseline supplied by the course, with modifications to work with text embeddings and classification.
    \subsection{Embedding Layer}
    Since the model operates on text data, an embedding layer is used to convert input tokens into dense vector representations.
    An important aspect of the embedding layer is the choice of vocabulary and tokenization strategy, which can significantly impact model performance.
    For this project two types of embeddings are used: word-level and subword-level embeddings.
    Word-level embeddings represent each unique word in the vocabulary as a vector, while subword-level embeddings break words into smaller units (e.g., using Byte Pair Encoding) to handle out-of-vocabulary words and capture morphological information.
    Since not all words in the dataset may be present in the embedding vocabulary, an \"unknown\" token is used to represent out-of-vocabulary words, allowing the model to still process them without crashing. 
    While initialising the embedding layer, the vocabulary size is set to a given value (e.g. 30000) and the weights are typically set to small random values, which allows the model to learn meaningful representations during training.
   \subsection{Multi-head Attention}
    Multi-head attention is the central mechanism that enables the model to attend to different parts of an input sequence in parallel.
    Instead of computing a single attention distribution, the mechanism uses multiple heads, each with its own learned projections, to capture different types of relationships in the data (for example, short- vs.\ long-range dependencies or syntactic vs.\ semantic information).
    The outputs of all heads are then concatenated and linearly transformed, yielding a richer and more expressive representation of the input sequence.

    \paragraph{Hyperparameters}
    \begin{itemize}
        \item Embedding size: 128
        \item Vocabulary size / encoded words: 30{,}000
        \item Number of attention heads: 4
        \item Dropout: 0.3
        \item Epochs: 50
    \end{itemize}

    In our experiments, this configuration achieves an accuracy of 0.82 on the training (own) dataset, while the accuracy on the held-out evaluation dataset is 0.38.
    This discrepency suggests that there is a missalignment between the training and evaluation datasets, which could be due to differences in data distribution, vocabulary, or other factors.
    To adress this issue a more in depth comparison between the two datasets is needed, which could include an analysis of the vocabulary overlap, distribution of classes, and other relevant features.
\section{Complementary Models Within the Project }

In addition to the other approaches explored in the project, this section presents three complementary models for the text classification task: a TF--IDF-based Feedforward Neural Network, a DistilBERT transformer, and a hybrid TF--IDF model combined with a Linear Support Vector Classifier. All three models were developed within the shared project pipeline, including common preprocessing steps, balanced sampling, and the same five-class classification objective.

The three models were chosen to represent different levels of complexity and different ways of encoding textual information. The first approach uses sparse lexical features and a simple neural classifier, providing a lightweight baseline. The second model moves to transformer-based contextual representations, allowing the system to capture semantic dependencies within the sentence. The third approach remains in the classical machine learning setting but strengthens it through a richer handcrafted representation that combines both word-level and character-level information. Together, these models contribute to the overall comparative analysis of the project by exploring the trade-off between computational efficiency, interpretability, and contextual understanding.

\subsection{TF--IDF + Feedforward Neural Network}

The first model is a classical text classification pipeline that combines TF--IDF feature extraction with a Feedforward Neural Network implemented in PyTorch. After selecting the relevant text and label columns, the data is cleaned, labels are mapped to the five target classes, and the dataset is split using an 80/20 stratified partition. The textual input is then transformed into a sparse TF--IDF representation with a maximum of 5000 features and English stopwords removed. In this way, each document is encoded as a fixed-dimensional vector capturing the relative importance of its words within the corpus.

These vectors are used as input to a fully connected neural architecture composed of two hidden layers. The first hidden layer projects the input representation to 256 neurons, and the second reduces it to 128 neurons. Both layers are followed by ReLU activation and dropout with rate 0.3, while the final output layer predicts one of the five classes. The model is trained using CrossEntropyLoss and the Adam optimizer with learning rate 0.001 for 5 epochs. This solution is computationally efficient and easy to train, making it a useful baseline inside the project. At the same time, its main limitation is that TF--IDF captures lexical salience but not true contextual meaning, which can reduce performance when class distinctions depend on semantic nuances rather than simple word frequency.

\begin{figure}[H]
    \centering
    \includegraphics[width=\textwidth]{tfidf_ffnn_pipeline.png}
    \caption{Pipeline of the TF--IDF + Feedforward Neural Network model.}
    \label{fig:tfidf_ffnn_pipeline}
\end{figure}

\subsection{DistilBERT Transformer}

The second model is based on DistilBERT, a compact transformer architecture derived from BERT. Unlike TF--IDF approaches, DistilBERT does not treat a document as an unordered set of words. Instead, it processes tokenized text through self-attention layers, producing contextual representations in which the meaning of each token depends on the surrounding sequence. This makes the model particularly suitable for a task such as source-family classification, where subtle contextual cues may help distinguish between human-written text and outputs generated by different AI families.

The training pipeline begins with the same label-cleaning and mapping logic used in the broader project. Labels are encoded numerically, and the data is split into training and validation sets using a stratified 80/20 partition. The pretrained checkpoint chosen for fine-tuning is \texttt{distilbert-base-uncased}. Texts are tokenized into subword units with a maximum sequence length of 256 tokens, and batching is handled dynamically.

For classification, \texttt{AutoModelForSequenceClassification} 
is used to place a task-specific classification head on top of the pretrained encoder. Fine-tuning is carried out through the Hugging Face Trainer API for 3 epochs, with learning rate $2 \times 10^{-5}$, batch size 16, weight decay 0.01, and early stopping to prevent unnecessary training once validation performance stops improving. Compared with the previous model, DistilBERT is computationally heavier, but it offers a much stronger ability to capture semantic and contextual information, making it one of the most relevant models in the overall comparison.

\begin{figure}[H]
    \centering
    \includegraphics[width=\textwidth]{distilbert_pipeline.png}
    \caption{Pipeline of the DistilBERT-based text classification model.}
    \label{fig:distilbert_pipeline}
\end{figure}

\subsection{Hybrid TF--IDF + LinearSVC}

The third model is a hybrid classical machine learning system that combines TF--IDF representations at both word level and character level with a Linear Support Vector Classifier. While it does not rely on deep contextual embeddings, it is designed to enrich the classical feature-based approach and better capture stylistic regularities that may distinguish among source families. This is particularly important in a project where the task is not only to separate human and AI text, but also to discriminate between outputs produced by different companies or model families.

The preprocessing stage is flexible and robust: the script detects valid text and label columns, normalizes label values, maps raw model names to the final classes when needed, and removes duplicated text-label pairs. The cleaned data is then split with an 80/20 stratified partition. The core of the method lies in the hybrid feature representation. At word level, TF--IDF is computed using n-grams from 1 to 3, with minimum document frequency equal to 2 and up to 50{,}000 features. At character level, TF--IDF is computed using character n-grams within word boundaries, with range from 3 to 6, minimum document frequency equal to 2, and up to 80{,}000 features. The two sparse matrices are concatenated into a single representation and used to train a \texttt{LinearSVC} model with balanced class weights and regularization parameter $C = 0.5$.

This model is an important complementary component of the project because it strengthens the traditional approach without requiring the cost of transformer fine-tuning. Word n-grams capture lexical and phrase-level patterns, while character n-grams are useful for stylistic details such as punctuation habits, suffixes, formatting tendencies, and token shapes. As a result, the model can be very competitive in authorship-related tasks. Its main limitation is that, despite the richer feature engineering, it still cannot model semantics and long-range dependencies as deeply as a transformer.

\subsection{Final Remarks}

Overall, these three models provide three different perspectives on the same classification problem. The TF--IDF + Feedforward Neural Network offers a simple neural baseline, DistilBERT introduces contextual language modeling, and the hybrid TF--IDF + LinearSVC provides a stronger sparse-feature alternative capable of capturing both lexical and stylistic information. Together, these models complement the broader project pipeline and contribute to the final comparison between efficient classical solutions and more advanced deep learning approaches.

\section{Conclusion}

In this project, we investigated the challenging task of classifying text samples according to their origin, distinguishing between human-written content and outputs generated by different large language models. To address the absence of a unified benchmark dataset, a substantial effort was devoted to constructing a diverse and balanced dataset by combining multiple publicly available sources. This step proved to be crucial, as the quality, diversity, and consistency of the dataset directly influence the effectiveness of the trained models.

Several modeling approaches were explored, ranging from classical machine learning techniques to modern transformer-based architectures. The custom multi-head self-attention transformer demonstrated the ability to learn meaningful representations from text; however, the observed discrepancy between training and evaluation performance highlighted issues related to domain shift and dataset misalignment. This finding emphasizes the importance of dataset quality, distribution matching, and generalization when dealing with real-world text classification problems. The TF--IDF-based neural network provided a computationally efficient baseline, while the hybrid TF--IDF + LinearSVC model showed strong performance in capturing stylistic and lexical patterns. The DistilBERT model, leveraging contextual embeddings, offered a more advanced understanding of semantic relationships and demonstrated the advantages of pretrained transformer architectures for this task.

\subsection*{Model Performance Summary}

\begin{itemize}
    \item Custom Transformer:
    \begin{itemize}
        \item Training Accuracy: 78.70
        \item Evaluation Accuracy: 36.67
    \end{itemize}
    
    \item TF--IDF + Feedforward Neural Network:
    \begin{itemize}
        \item Training Accuracy: 80.01
        \item Evaluation Accuracy: 38.04
    \end{itemize}
    
    \item DistilBERT:
    \begin{itemize}
        \item Training Accuracy: 92.33
        \item Evaluation Accuracy: 44.00
    \end{itemize}
    
    \item Hybrid TF--IDF + LinearSVC:
    \begin{itemize}
        \item Training Accuracy: 94.58
        \item Evaluation Accuracy: 44.00
    \end{itemize}
\end{itemize}

Overall, the results suggest that distinguishing between human and machine-generated text—and even between different model families—is feasible, but remains a complex problem, particularly when models are trained and evaluated on data from different distributions. Future work could focus on improving dataset alignment, incorporating more diverse and up-to-date model outputs, and exploring advanced techniques such as domain adaptation, ensemble methods, or larger pretrained models. Additionally, deeper analysis of linguistic and stylistic features could provide further insights into what differentiates human writing from machine-generated text.

In conclusion, this project demonstrates both the potential and the limitations of current approaches to text origin classification, and highlights the need for continued research as language models evolve and become increasingly sophisticated.
\end{document}
